This paper is the companion of the first paper of the series, \emph{Matrix Transposition and the Tape Coding
Conjecture}. That paper reduced the question whether a multitape Turing machine can transpose a $d\times d$ matrix of
$\lceil\log_2d\rceil$-bit entries in linear time to a conjecture about network codes on the time--space diagrams of
tapes. Here the implication runs the other way: a conjecture about arrays alone, with no machine, no time and no
network in its statement, implies the Tape Coding Conjecture.

A \emph{whole-row code} of window~$w$ is a set of $d\times d$ arrays in which every column is a function of at most
$w$ rows. The symmetric arrays are such a code, with $w=1$ and about half of all entries free. The \emph{Matrix Coding
Conjecture} says that there is a $\delta>0$ such that, for every fixed~$w$ and all large~$d$, a whole-row code of
window~$w$ has at most $2^{(1-\delta)L}$ elements, where $L$ is the number of bits of an array. We prove that the
Matrix Coding Conjecture implies the Tape Coding Conjecture in its robust form, and therefore that no multitape Turing
machine, oblivious or not, would transpose in linear time. A weak form, in which the deficit $\delta$ may shrink slowly
as the window grows and only one size $d$ beyond every threshold is needed, already suffices. For oblivious machines the
reduction is direct: after a cheap cut of the run, the bits on which the outputs of almost every row of the transpose
depend lie in boundedly many rows of the input.

The conjecture is open, and the rest of the paper is about it. Every whole-row code contains a large code in which rows
are also functions of columns, on a symmetric relation. Every bound obtained from an acyclic set of entries is the
bound of an order of the rows, and such bounds, with or without transported copies of the recovery rules, cannot give
a deficit that is independent of the window. Entropy inequalities can go further: on the incidence relation of the
Fano plane the best acyclic bound is $41$ entries of $49$ and the Shannon value is $119/3$, and the gap persists, at more
than one percent of all entries, on every lift whose helper rows respect the incidences of the plane and which has a
mixing property. Codes that are invariant under the shifts of a
pairing of the entries, with $b$ pairs omitted, have rate at most $\tfrac12+b/d^2+o(1)$ over every finite abelian group
and for every relation of window~$w$, with nonlinear recovery allowed.

The theorems above are formalized in Lean~4, with the standard axioms only, against the frozen statement of the first
paper. Proved in the text and not formalized are two existence statements (relations of bounded window with the mixing
property, and mixing lifts of the Fano relation), the code that attains the Fano optimum, and one further result: the
exclusion of oblivious runs whose time diagram a cheap cut splits into pieces with few outputs.
